% Ad Atticum 11.18 --- headnote
% ad-atticum-11-18, 20 June 47 BC, Brundisium

\textit{Cicero to Atticus, written from Brundisium on
the twelfth day before the Kalends of Quintilis (July)
47 BC --- 20 June (the manuscript dateline:
\textit{Scr.\ Brundisi xii K.\ Quint.\ a.\ 707 (47)}).
\textit{Ille} is Caesar, held up in Alexandria beyond
what Cicero had hoped: there is no rumour of his
departure, and the report is that the king's war is
detaining him severely. Cicero, who had thought of
sending his son back, abandons the plan and begs
Atticus to get him out of Brundisium --- any punishment
is lighter than staying on. He has written to the same
effect to Antony, Balbus, and Oppius. If the next
campaign falls in Italy, or if Caesar uses his fleets,
Brundisium is the worst possible place for a Pompeian
of doubtful standing to be sitting in; one of those
two contingencies, he reckons, is certain to come
about, and perhaps both.}

\textit{In the second section Atticus's earlier
report of Oppius's conversation comes into focus:
Cicero has understood from it how bitter the Caesarian
inner circle is against him, and asks Atticus to soften
that anger. The condition he describes is the bottom
of the book --- ``I expect nothing now but misery, but
nothing more ruined than what I am now in can possibly
happen.'' What he wants done is precise: Atticus is to
talk to Antony and the others, extricate the matter
as best he can, and write back on every point as soon
as possible. The signed dateline is preserved at the
foot of the letter.}
